\section{幂修正}\label{sect_powercorrections}%
%

式 \eqref{BorelIntegral} 中积分路径上的奇点意味着微扰论不完备，级数之和没有良好定义。
%Following ~\cite{Beneke:1998ui,Beneke:2000kc}
按惯例 ~\cite{Beneke:1998ui}，相应的模糊性可估计为
\begin{align}
\delta H(w_0)  = -\pi \frac{1}{\beta_0}  e^{-w_0/(\beta_0 a_s)} \res_{w=w_0}\big[B[H](w)\big]\,,
\end{align}
其中 $w_0$ 是奇点位置，$\displaystyle{\res_{w=w_0}\big[B[H](w)\big]}$ 是相应留数。注意 $ e^{-w_0/(\beta_0 a_s)} = (\Lambda^2/\mu^2)^{w_0}$。按照标准逻辑~\cite{Beneke:1998ui,Beneke:2000kc}，我们假定这一模糊性必须通过添加同量级的非微扰修正来抵消。

%
\subsection{Ioffe 时间准分布}\label{sect_qITD}%
%

考察 $\delta H(1)$，我们得到以``Ioffe 时间''
\begin{align}
 \tau = (p\cdot v) z
\end{align}
为变量的 qITD 领头幂修正：
\begin{align}
\label{t4-qITD}
  \mathcal{I}^\parallel(\tau)
&= I(\tau) +
  \kappa (v^2\! z^2\!\Lambda^2)\! \int_0^1\!d\alpha \, (\alpha + \bar\alpha \ln \bar\alpha)I(\alpha \tau )\,,
\notag\\
  \mathcal{I}^\perp(\tau)
%(z^2v^2, pvz, \mu)
&= I(\tau) +
  \kappa (v^2\! z^2\!\Lambda^2) \!\int_0^1\!d\alpha \, (\alpha + \bar\alpha \ln \bar\alpha + \alpha\bar\alpha)I(\alpha \tau )\,,
\end{align}
其中 $\kappa$ 是量级为一的实数。\eqref{w=1} 对应的 renormalon 模糊性给出
\begin{align}
  |\kappa| = -\pi \frac{1}{\beta_0}\left(- \frac{ e^{5/3}}{4 e^{-2\gamma_E}}\right) = \frac{\pi e^{5/3+2\gamma_E}}{4 \beta_0} \simeq 1.5\,,
\end{align}
但这个数只有参考意义。另一种做法是取 $\kappa=1$，并把 $\Lambda$ 看作 $\Lambda_{\rm QCD}$ 量级的某个非微扰参数，它决定幂修正的整体归一化且在此方法中无法确定。注意修正的符号同样未被确定。

对于式 \eqref{n-qITD} 定义的归一化 qITD，我们得到的则是
\begin{align}
\label{t4-nqITD}
  \mathbf{I}^\parallel(\tau)
%(z^2v^2, pvz, \mu)
&= I(\tau) +
  \kappa (v^2\! z^2\!\Lambda^2) \!\int_0^1\!d\alpha \, [\alpha + \bar\alpha \ln \bar\alpha]_+I(\alpha \tau )\,,
\notag\\
  \mathbf{I}^\perp(\tau)
%(z^2v^2, pvz, \mu)
&= I(\tau) +
  \kappa (v^2\! z^2\!\Lambda^2) \!\int_0^1\!\!d\alpha \, [\alpha + \bar\alpha \ln \bar\alpha + \alpha\bar\alpha]_+I(\alpha \tau )\,,
\end{align}
其中``加号''分布按通常方式定义：
\begin{align}
  [f(\alpha)]_+ = f(\alpha) - \delta(\bar\alpha)\int_0^1 d\beta\, f(\beta)\,.
\end{align}
按真空关联函数归一化的 qITD $\widehat{\mathbf{I}}^{\parallel(\perp)}$~\eqref{hat-qITD} 的领头幂修正与未归一化分布 \eqref{t4-qITD} 的相同。
式~\eqref{t4-qITD} 与 \eqref{t4-nqITD} 是以下分析的出发点。

为了直观展示幂修正相对于领头扭度结果对``Ioffe 时间''$\tau$ 的函数依赖，我们写成
\begin{align}
    \cI &= I(\tau)\Big\{ 1 + \kappa (v^2 z^2\Lambda^2) \cR_{\cI}(\tau)\Big\},
\notag\\
    \bI &= I(\tau)\Big\{ 1 + \kappa (v^2 z^2\Lambda^2) \bR_{\cI}(\tau)\Big\},
\end{align}
其中
\begin{align}
\cR_{\cI}^\parallel(\tau) & = \frac{1}{I(\tau)} \int_0^1\!d\alpha \, 
(\alpha + \bar\alpha \ln \bar\alpha)I(\alpha \tau)\,,
\notag\\  
\cR_{\cI}^\perp(\tau)
& = \frac{1}{I(\tau)} \int_0^1\!d\alpha \, 
(\alpha + \bar\alpha \ln \bar\alpha + \alpha\bar\alpha)I(\alpha \tau)\,.
\end{align}
对归一化 qITD 可得
\begin{align}
  \bR_{\cI}^\parallel(\tau) &= \cR_{\cI}^\parallel(\tau) - \frac14\,, 
\qquad
  \bR_{\cI}^\perp(\tau) &= \cR_{\cI}^\perp(\tau) - \frac5{12}\,, 
\end{align}
且显然有
\begin{eqnarray}
\widehat{\bR}_{\cI}^{\parallel(\perp)}(\tau) = \cR_{\cI}^{\parallel(\perp)}(\tau),
\end{eqnarray}
故不再单独讨论。

一般而言，qITD $\cI(\tau)$ 及高扭度系数 $\cR_{\cI}(\tau)$、$\bR_{\cI}(\tau)$ 是复函数，但其虚部看起来很小。对于价夸克分布的一个简单模型
\begin{align}
\label{q-model}
  q(x) = x^{-1/2}(1-x)^3
\end{align}
其实部 $\Real\cR_{\cI}(\tau)$ 与 $\Real\bR_{\cI}(\tau)$ 绘于图~\ref{fig:R-qITD}。
%
\begin{figure}[t]
\centering
\includegraphics[width=6.0cm]{BVZplots/RImodel}~
\caption{式~\text{\eqref{q-model}} 的简单夸克 PDF 模型下
$\cR_{\cI}(\tau)$（红色粗线）与 $\bR_{\cI}(\tau)$（黑色细线）的实部。
实线为 $\Real R^\parallel$，虚线为 $\Real R^\perp$。}
\label{fig:R-qITD}
\end{figure}
%
结果表明，``横向'' qITD 的幂修正大约是``纵向'' qITD 的四倍。
两种情形下，$R$ 函数在大 Ioffe 时间 $\tau \ge 10$ 处趋于平缓，然而现今的格点计算几乎难以达到这样的 $\tau$。按零质子动量归一化相当于减去 $\tau=0$ 处的取值，
因此在 $z\to 0$ 时按构造不存在幂修正。

%
\subsection{部分子准分布}\label{sect_qPDF}%
%

对上述 qITD 结果作 Fourier 变换即得 qPDF：
%
\begin{figure*}[t]
\centering
 \includegraphics[width=6.0cm]{BVZplots/RQparallel}
 ~\includegraphics[width=6.2cm]{BVZplots/RQtrans-par}
 %~\includegraphics[width=6.0cm]{BMJplots/fit-t4-qPDF}
\caption{$\cR^\parallel_\cQ(x)$（左图）与 $\cR^\perp_\cQ(x) - \cR^\parallel_\cQ(x)$（右图）：
MSTW 价 u 夸克（黑色长虚线）、d 夸克（蓝色短虚线）（均在 2 GeV）以及
式~\eqref{q-model} 的模型（红色实线）。}
\label{fig:calR}
\end{figure*}
%
\begin{align}
\label{t4-Qparallel}
 \cQ^\parallel(x,p) &= 
q(x) - \frac{\kappa v^2\! \Lambda^2}{(pv)^2}\Big(\frac{d}{dx} \Big)^2 \int_{|x|}^1\!\frac{dy}{y} \, (y + \bar y \ln \bar y) q(\tfrac{x}{y})
%\notag\\&= q(x) - \frac{\kappa v^2 \!\Lambda^2}{x^2(pv)^2} \int_{|x|}^1\!\frac{dy}{y} \Big[\frac{y^2}{[1-y]_+} + 2\delta(\bar y) 
%\notag\\&\qquad 
%+\delta'(\bar y)\Big]q(\tfrac{x}{y})
\notag\\&= q(x) - \frac{\kappa v^2 \!\Lambda^2}{x^2(pv)^2} 
\biggl\{
\int_{|x|}^1\!\frac{dy}{y} \frac{y^2}{[1-y]_+} q(\tfrac{x}{y})
\notag\\&\qquad
 + q(x) - |x|q'(x) \biggr\}
\end{align}
以及
\begin{align}
\label{t4-Qperp}
 \cQ^\perp(x,p) &= 
q(x) - \frac{\kappa v^2\! \Lambda^2}{(pv)^2}\Big(\frac{d}{dx} \Big)^2\!\!\! \int_{|x|}^1\!\frac{dy}{y}  (y + \bar y \ln \bar y + y\bar y) q(\tfrac{x}{y})
%\notag\\&=
% q(x) - \frac{\kappa v^2 \!\Lambda^2}{x^2(pv)^2}
% \int_{|x|}^1\!\frac{dy}{y} 
%\Big[\frac{y^2}{[1-y]_+} + 3\delta(\bar y) 
%\notag\\&\qquad 
%+\delta'(\bar y)-2 y^2\Big]
%q(\tfrac{x}{y})\,.
\notag\\&= q(x) - \frac{\kappa v^2 \!\Lambda^2}{x^2(pv)^2} 
\biggl\{
\int_{|x|}^1\!\frac{dy}{y} \left[\frac{y^2}{[1-y]_+} - 2 y^2\right] q(\tfrac{x}{y})
\notag\\&\qquad
 + 2 q(x) - |x|q'(x) \biggr\}
\end{align}
以下设 $x>0$，则这些表达式可以改写为
\begin{align}
\label{cR}
   \cQ^{\parallel(\perp)}(x,p) &=  q(x) \biggl\{1 - \frac{v^2 \Lambda^2}{ x^2\bar x  (pv)^2}\cR^{\parallel(\perp)}_\cQ(x) \biggr\}
\end{align}
其中
\begin{align}
 R^\parallel_Q(x) &= 
\frac{\bar x }{q(x)} \biggl\{ \int_{x}^1\!\frac{dy}{1-y} \Big[y q(\tfrac{x}{y}) - q(x)\Big] + q(x)-x q'(x)\biggr\}, 
\notag\\
 R^\perp_Q(x) &= 
\frac{\bar x }{q(x)} \biggl\{ \int_{x}^1\!\frac{dy}{1-y} \Big[y (2y-1) q(\tfrac{x}{y}) - q(x)\Big] + 2q(x)  
\notag\\&\qquad - x q'(x) \biggr\}.
\end{align}
%
注意我们为幂修正提取出前置因子 $1/(x^2\bar x)$，预期它在 Bjorken 变量很小（$x\to 0$）和很大（$x\to 1$）的区域分别被放大为 $1/x^2$ 与 $1/(1-x)$。

归一化 QPDF $\bQ^{\parallel(\perp)}(x,p)$ 由把式~\eqref{t4-Qparallel} 与 \eqref{t4-Qperp} 第一行中的核替换为加号分布而得，参见 \eqref{t4-nqITD}。把结果写成
\begin{align}
\label{bR}
   \bQ^{\parallel(\perp)}(x,p) &=  \hat q(x) \biggl\{1 - \frac{v^2 \Lambda^2}{ x^2\bar x  (pv)^2}\bR^{\parallel(\perp)}_\cQ(x) \biggr\}, 
\end{align}
其中 $\hat q(x)$ 是积分为一的归一化夸克 PDF，于是得到
\begin{align}
 \bR^\parallel_\cQ(x) &= \cR^\parallel_\cQ - \frac14 \frac{\bar x x^2 }{q(x)}  q''(x)\,, %= \cR^\parallel_Q + \delta R^\parallel_\cQ\,,
\notag\\
 \bR^\perp_\cQ(x) &= \cR^\perp_\cQ - \frac5{12} \frac{\bar x x^2 }{q(x)}  q''(x)\,.    % = \cR^\perp_Q + \delta R^\perp_\cQ\,.
\end{align}
注意附加项正比于夸克 PDF 的二阶导数，因此在 $x\to 1$ 处以 $1/(1-x)^2$ 增强。
如前所述，对真空关联函数的归一化不影响 $1/p^2$ 幂修正，因此
$\widehat{\bR}_{\cQ}^{\parallel(\perp)}(\tau) = \cR_{\cQ}^{\parallel(\perp)}(\tau)$，无需单独讨论该情形。

数值研究采用了 2 GeV 尺度下的 MSTW NLO 价 $u$、$d$ 夸克分布 \cite{Martin:2009iq} 以及式 \eqref{q-model} 的简单模型。结果显示，（未归一化的）``纵向''与``横向''qPDF 的幂修正大小相近，因此我们在图~\ref{fig:calR} 左图画出 $\cR^\parallel_\cQ(x)$，右图画出差值
$\cR^\perp_\cQ(x) - \cR^\parallel_\cQ(x)$。

对所有夸克 PDF 模型，$\cR_\cQ(x)$ 的 $x$ 依赖都类似：幂修正在 $x\to 0$ 时很小
（但不为零：$\cR^\parallel_\cQ(0)\sim 0.2$、$\cR^\perp_\cQ(0)\sim 0.4$），并随 $x$（近乎线性地）陡增。对 $u$ 夸克，结果与式 \eqref{q-model} 的简单模型非常接近；而对 $d$ 夸克，幂修正约大一倍。注意差值 $\cR^\perp_\cQ(x) - \cR^\parallel_\cQ(x)$ 对 PDF 模型的依赖非常微弱。

用零动量处归一化的 qITD 构造 qPDF 会产生很大影响。
图~\ref{fig:bRparallel} 展示了这一点：上排是``纵向''情形的结果，下排是``纵向''与``横向''分布之差。
%
\begin{figure*}[tbp]
\centering
 \includegraphics[width=5.0cm]{BVZplots/normRQu}~
 \includegraphics[width=5.0cm]{BVZplots/normRQd}~
 \includegraphics[width=5.0cm]{BVZplots/normRQm}
\\
 \includegraphics[width=5.0cm]{BVZplots/deltaR-RQu}~
 \includegraphics[width=5.0cm]{BVZplots/deltaR-RQd}~
 \includegraphics[width=5.0cm]{BVZplots/deltaR-RQm}
\caption{上排：MSTW 价 u 夸克（左）、d 夸克（中）（均为 2 GeV）以及
式~\eqref{q-model} 模型（右）的
$\cR^\parallel_\cQ$（红色实线）与 $\bR^\parallel_\cQ$（黑色虚线）。
下排：$\cR^{\perp}_\cQ-\cR^{\parallel}_\cQ $ 与 $\bR^{\perp}_\cQ-\bR^{\parallel}_\cQ $ 的对应结果。
注意分栏坐标：$0<x<0.6$ 与 $0.6<x<1$ 两段在纵轴上使用了不同标度。}
\label{fig:bRparallel}
\end{figure*}
%
三幅子图（从左到右）分别对应 MSTW 价 u 夸克、d 夸克以及式 \eqref{q-model} 的简单模型。红色实线表示 $\cR_\cQ(x)$（与图~\ref{fig:calR} 相同），归一化 qPDF 的结果 $\bR_\cQ(x)$ 用黑色虚线画出。

可以看到，按零动量 qITD 归一化显著减小了中等 $x \lesssim 0.5$ 区域的幂修正，但代价是在更大 $x$ 处急剧增大。因此这一归一化步骤不适合研究 PDF 的大-$x$ 行为，但显然能在不太大的 $x$ 区域使幂修正最小化。

%
\subsection{部分子赝分布}\label{sect_pPDF}%
%

\begin{figure}[t]
\centering
 \includegraphics[width=5.0cm]{BVZplots/pseudoR}~
% \includegraphics[width=5.0cm]{BMJplots/t4-pPDF}~
\caption{\label{fig:pPDF} pPDF 的幂修正 $\mathcal{R}_\cP$ \eqref{t4-pPDF}：
MSTW 价 u 夸克（黑色长虚线）、d 夸克（蓝色短虚线）（均为 2 GeV），以及简单模型
$q(x) = x^{-1/2} (1-x)^3 $（红色实线）。}
\end{figure}

pPDF 的幂修正可以直接由``横向''qITD 的相应表达式 \eqref{t4-qITD}、\eqref{t4-nqITD} 得到。
把结果写成
\begin{align}
\mathcal{P}(x,z,\mu) &= q(x)\Big\{1 + (v^2\! z^2\!\Lambda^2)  \theta(|x|<1) \mathcal{R}_\cP(x)\Big\},
\end{align}
并对按零动量归一化的 pPDF 即 $\mathbf{P}(x,z)$ 作类似处理，可得
\begin{align}
\label{t4-pPDF}
\mathcal{R}_\cP(x) &= \frac{1}{q(x)} \int_{|x|}^1\!\frac{dy}{y} \, (y + \bar y \ln \bar y +  y \bar y ) q(\tfrac{x}{y})\,,
\notag\\
\mathbf{R}_\cP(x) &= \frac{1}{q(x)} \int_{|x|}^1\!\frac{dy}{y} \, [y + \bar y \ln \bar y +  y \bar y ]_+ q(\tfrac{x}{y})
\notag\\&=
\mathcal{R}_\cP(x) - 5/12 \,.
\end{align}
数值结果示于图~\ref{fig:pPDF}。

注意对所有所考虑的价夸克 PDF 模型，$\mathcal{R}_\cP(x)$ 都十分相似，并在 $x\to 1$ 处减小。事实上，容易看出在该极限下 $\mathcal{R}_\cP(x) = \mathcal{O}(1-x)$，类似于靶质量修正（式~\eqref{target-pPDF}）。一旦把 pPDF 按零动量归一化（这会附加一个负常数），这一压低便不复存在；但如果按同一算符的真空期望值归一化，则压低得以保持。

%
